\capitulo{8}{Evaluación del modelo y discusión}

En este apartado se hablará de como ha sido evaluado el modelo por los tutores del proyecto con el fin de verificar si cumple o no con los requisitos establecidos en el apartado \ref{ap:analisis_requisitos}

Este apartado esta basado en los resultados encuesta realizada a los expertos que han probado el modelo, que se encuentra en el apendice \ref{ap:encuesta_expertos}

\section{Evaluación del modelo}

\subsection{Resolución de ejercicios}

En esta prueba se ha evaluado la capacidad del modelo para resolver ejercicios, es decir que se le pase un enunciado y genere un código que se asemege lo máximo posible a lo que el usuario introduce en lenguaje natural.

Tras diferentes pruebas realizadas por los profesores sobre el modelo estos le han dado la siguiente calificacion.

\begin{table}[htb]
    \centering
    \begin{tabular}{| >{\arraybackslash}m{4,5cm}| >{\arraybackslash}m{4,5cm} | >{\arraybackslash}m{4,5cm}|} \hline
    \vspace{0.5em}  Muy incorrecto (incluye errores de mala praxis) \vspace{0.5em} &  Incluye alguna deficiencia menor de eficiencia (sin mala praxis) & Totalmente correcto (óptimo con respecto a eficiencia y uso de estructuras de control) \\ \hline
     & \multicolumn{1}{|c|}{X} & \\ \hline 
    \end{tabular}\caption{Prueba resolución de ejercicios modelo}
    \label{tabla:evaluacion-ej1}
\end{table}

En esta prueba el modelo se ha comportado correctamente, es cierto que a la hora de generar código mete librerías java que en un programa normal serían una buena práctica, pero dado el contexto de la asignatura que como se ha visto en el análisis de requisitos presenta restricciones, el modelo no consigue ceñirse estrictamente al codigo permitido en la asignatura de FPROG.

\subsection{Generación de un enunciado en lenguaje natural}
\label{generacion_enunciado}
En esta prueba se pretende que el modelo genere un enunciado en lenguaje natural, con el fin de que los alumnos puedan obtener diferentes problemas con los que prácticar.

En esta tarea se consiguieron generar diferentes enunciados, uno de ellos siendo el siguiente

\imagen{conversacionModelo1}{Prompt realizado por el usuario y respuesta del modelo \cite{openwebui}}

Se consiguieron generar mas enunciados, aun así es cierto que las pruebas han indicado que cuesta obtener este tipo de promt sin que incluya nada de código, lo que puede dificultar o frustrar a los alumnos en el uso del modelo.

\subsection{Evaluación de codigo correcto por parte del modelo}

En esta prueba se evalua si al pasarle un código java al modelo este es capaz de evaluarlo según los requisitos de la asignatura de FPROG, indicando  que es correcto.

En las pruebas realizadas por los expertos han sido satisfactorias estos han indicado la siguiente calificación

\begin{table}[htb]
    \centering
    \begin{tabular}{| >{\arraybackslash}m{4,5cm}| >{\arraybackslash}m{4,5cm} | >{\arraybackslash}m{4,5cm}|} \hline
    \vspace{0.5em}  Evaluación incorrecta (detecta errores que no existen) \vspace{0.5em} &  Evaluación correcta (sin mensajes explicativos) & Evaluación totalmente correcta (incluye mensajes explicativos) \\ \hline
     &  & \multicolumn{1}{|c|}{X} \\ \hline 
    \end{tabular}\caption{Prueba evaluación de codigo correcto por parte del modelo}
    \label{tabla:evaluacion-ej1}
\end{table}

\subsection{Evaluación de código con error de mala praxis}

En esta prueba se evalua si al pasarle un código en lenguaje Java que contenga alguno de los errores que se encuentran en el apartado \ref{ap:analisis_requisitos}, el modelo debe ser capaz de informar en su respuesta al promt de los errores dde malas praxis detectados, asi como generar un código que arregle esos errores.

En los casos probado por los expertos se ha dado la siguiente valoración

\begin{table}[htb]
    \centering
    \begin{tabular}{| >{\arraybackslash}m{4,5cm}| >{\arraybackslash}m{4,5cm} | >{\arraybackslash}m{4,5cm}|} \hline
    \vspace{0.5em}  Evaluación incorrecta (detecta errores que no existen) \vspace{0.5em} &  Evaluación correcta (sin mensajes explicativos) & Evaluación totalmente correcta (incluye mensajes explicativos) \\ \hline
     &  & \multicolumn{1}{|c|}{X} \\ \hline 
    \end{tabular}\caption{Prueba evaluación de código con error de mala praxis}
    \label{tabla:evaluacion-ej1}
\end{table}

En las pruebas realizadas por los expertos se ha conseguido el objetivo satisfactoriamente con diferentes problemas de programación vistos en la asignatura.

\subsection{Evaluación generación de código con iteraciones}

En esta evaluación se solicita que el modelo genere un enunciado que trate las iteraciones, en este caso como sucede en la evaluacion \ref{generacion_enunciado}, cuesta que el modelo genere una respuesta adecuada a la petición del usuario.

\imagen{conversacionModelo2}{Prompt realizado por el usuario y respuesta del modelo \cite{openwebui}}

\subsection{Evaluación respuesta ejercicio corregido}

En esta prueba consiste en suministrar el código corregido y ver si lo indica como que es un código correcto.

En los casos probado por los expertos se ha dado la siguiente valoración

\begin{table}[htb]
    \centering
    \begin{tabular}{| >{\arraybackslash}m{4,5cm}| >{\arraybackslash}m{4,5cm} | >{\arraybackslash}m{4,5cm}|} \hline
    \vspace{0.5em}  Evaluación incorrecta (detecta errores que no existen) \vspace{0.5em} &  Evaluación correcta (sin mensajes explicativos) & Evaluación totalmente correcta (incluye mensajes explicativos) \\ \hline
     &  & \multicolumn{1}{|c|}{X} \\ \hline 
    \end{tabular}\caption{Prueba evaluación de respuesta ejercicio corregido}
    \label{tabla:evaluacion-ej1}
\end{table}

\subsection{Evalación de detección de punto de rotura de la estructuras de control}

En esta prueba se valora si el modelo es capaz de detectar el uso \textit{break} u otras instrucciones que provoquen una modificación intencionada de una estructuras de control, ademas de generar un código que elimine esa problemática.

En los casos probado por los expertos se ha dado la siguiente valoración

\begin{table}[htb]
    \centering
    \begin{tabular}{| >{\arraybackslash}m{4,5cm}| >{\arraybackslash}m{4,5cm} | >{\arraybackslash}m{4,5cm}|} \hline
    \vspace{0.5em}  Evaluación incorrecta (detecta errores que no existen) \vspace{0.5em} &  Evaluación correcta (sin mensajes explicativos) & Evaluación totalmente correcta (incluye mensajes explicativos) \\ \hline
     &  & \multicolumn{1}{|c|}{X} \\ \hline 
    \end{tabular}\caption{Prueba evaluación de detección de punto de rotura de la estructuras de control}
    \label{tabla:evaluacion-ej1}
\end{table}

\subsection{Evaluación de detección de bucles infinitos}

En esta prueba se valora si el modelo es capaz de detectar en el código bucles sean infinitos, inidcando al usuario el error y la correción del mismo.

En los casos probado por los expertos se ha dado la siguiente valoración

\begin{table}[htb]
    \centering
    \begin{tabular}{| >{\arraybackslash}m{4,5cm}| >{\arraybackslash}m{4,5cm} | >{\arraybackslash}m{4,5cm}|} \hline
    \vspace{0.5em}  Evaluación incorrecta (detecta errores que no existen) \vspace{0.5em} &  Evaluación correcta (sin mensajes explicativos) & Evaluación totalmente correcta (incluye mensajes explicativos) \\ \hline
     &  & \multicolumn{1}{|c|}{X} \\ \hline 
    \end{tabular}\caption{Prueba evaluación de detección de bucles infinitos}
    \label{tabla:evaluacion-ej1}
\end{table}

\section{Discusión rendimiento del modelo}

Tras la evaluacion de los expertos y las pruebas realizadas podemos concluir que el modelo realiza correctamente las labores de generación de código asi como de interpretación basandose en el documento de malas praxis como guia para diferenciar entre un código correcto u incorrecto según los criterios de la asignatura.

Aunque presenta problemas que debido a la extensión en la fase de investigación se tendrán en cuenta en las líneas de trabajo futuras \ref{ap:lineas de trabajo futuras}, las principales tareas en las que presenta dificultades el modelo son:

\begin{itemize}
    \item Generación de enunciados en lenguaje natural que no incluyan código: tras las pruebas se ha visto que el modelo necesita de múltiples prompts para que genere un enunciado que no contenga código, por lo que el modelo necesita mas ejemplos de este tipo en el conjunto de datos con el fin de que sea capaz de inferir este tipo de respuestas.
    
    \item Falsos positivos en el uso de mala praxis en el código: en algunas ocasiones el modelo detecta codigo que identifica como mala praxis cuando en realidad no lo es, lo que podría llevar a la confusión de los alumnos. Aunque la mayoria de veces lo realiza correctamente.
\end{itemize}




